\subsection{约当-赫尔德定理}

定义9. 带算子群G的一个合成列是有限序列 \((\mathbf{G}_i)_{0\leqslant i\leqslant n}\) G的稳定子群，其中 \(\mathbf{G}_0 = \mathbf{G}\) 和 \(\mathbf{G}_n = \{e\}\) 且使得 \(\mathbf{G}_{i + 1}\) 是由 \(\mathbf{G}_i\) 对于 \(0\leqslant i\leqslant n - 1\) 。这些商 \(\mathrm{G_i / G_{i + 1}}\) 被称为该序列的商。一个合成序列 \(\Sigma^{\prime}\) 被称为比一个合成序列更细 \(\Sigma\) 中正规，若 \(\Sigma\) 是从中取出的一个序列 \(\Sigma^{\prime}\) .

上定义了一个处处有定义的合成律。若 \((\mathbf{G}_i)_{0\leqslant i\leqslant n}\) 和 \((\mathbf{H}_j)_{0\leqslant j\leqslant m}\) 分别是两个带算子群的合成群列 \(\mathbf{G}\) 和 \(\mathbf{H}\) ，则称它们是等价的，如果 \(m = n\) 且存在一个排列 \(\phi\) 区间的 \([0,n - 1]\) 于 \(\mathbf{N}\) ，使得带算子的群 \(\mathbf{G}_i / \mathbf{G}_{i + 1}\) 和 \(\mathrm{H}_{\phi (i)} / \mathrm{H}_{\phi (i) + 1}\) 对所有 \(i\) .

注意，一般来说，取自合成列的序列 \((\mathbf{G}_{i})\) 不是合成列，因为对于 \(j > i + 1\) , \(G_{j}\) 一般不构成 \(G_{i}\) .

定理5（施赖埃尔）。给定两个合成列 \(\Sigma_{1}, \Sigma_{2}\) 带算子群G的两个合成列，存在两个等价的合成列 \(\Sigma_{1}^{\prime}, \Sigma_{2}^{\prime}\) ，它们分别加细于 \(\Sigma_{1}\) 和 \(\Sigma_{2}\) .

设 \(\Sigma_{1} = (\mathrm{H}_{i})_{0\leqslant i\leqslant n}\) 和 \(\Sigma_{2} = (\mathrm{K}_{j})_{0\leqslant j\leqslant p}\) 设为两个给定的合成列，分别具有 \(n + 1\) 和 \(p + 1\) 项；我们将看到合成列 \(\Sigma_1^\prime\) 可以由插入形成 \(p - 1\) 子群 \(\mathbf{H}_{i,j}^{\prime}\) \((1 \leqslant j \leqslant p - 1)\) 介于 \(H_{i}\) 和 \(H_{i+1}\) 对于 \(0 \leqslant i \leqslant n - 1\) 与该序列之间 \(\Sigma_{2}'\) 通过插入 \(n - 1\) 子群 \(K_{j,i}'\) ( \(1 \leqslant i \leqslant n - 1\) ) 于 \(K_{j}\) 和 \(K_{j+1}\) 对于 \(0 \leqslant j \leqslant p - 1\) 之间；从而将得到 G 的两列 \(pn + 1\) 稳定子群；通过适当选择所插入的稳定子群，我们将证明这些列是等价的合成列。

为此，注意 \(\mathbf{H}_i \cap \mathbf{K}_j\) 是 \(\mathbf{H}_i\) 以及 \(\mathbf{K}_j\) ，因此（定理4） \(\mathbf{H}_{i+1} \cdot (\mathbf{H}_i \cap \mathbf{K}_j)\) 是 \(\mathbf{H}_i\) ，包含 \(\mathbf{H}_{i+1}\) 和 \(\mathbf{K}_{j+1} \cdot (\mathbf{H}_i \cap \mathbf{K}_j)\) 是 \(\mathbf{K}_j\) ，包含 \(\mathbf{K}_{j+1}\) 。如果我们写 \(\mathbf{H}_{i,j}' = \mathbf{H}_{i+1} \cdot (\mathbf{H}_i \cap \mathbf{K}_j)\) 和 \(\mathbf{K}_{j,i}' = \mathbf{K}_{j+1} \cdot (\mathbf{H}_i \cap \mathbf{K}_j)\) , \(\mathbf{H}_{i,j+1}'\) 是 \(\mathbf{H}_{i,j}'\) ( \(0 \leqslant j \leqslant p-1\) ) 且 \(\mathbf{K}_{j,i+1}'\) 是 \(\mathbf{K}_{j,i}'\) ( \(0 \leqslant i \leqslant n-1\) ）。此外 \(\mathbf{H}_{i,0}' = \mathbf{H}_i\) , \(\mathbf{H}_{i,p}' = \mathbf{H}_{i+1}\) , \(\mathbf{K}_{j,0}' = \mathbf{K}_j\) 和 \(\mathbf{K}_{j,p}' = \mathbf{K}_{j+1}\) 为了证明该定理，只需证明 \(\mathbf{H}_{i,j+1}'\) （相应地 \(\mathbf{K}_{j,i+1}'\) ) 是 \(\mathbf{H}_{i,j}'\) （相应地 \(\mathbf{K}_{j,i}'\) ) 的一个正规稳定子群，并且商群 \(\mathbf{H}_{i,j}'/\mathbf{H}_{i,j+1}'\) 和 \(\mathbf{K}_{j,i}'/\mathbf{K}_{j,i+1}'\) 是同构的（ \(0 \leqslant i \leqslant n-1\) , \(0 \leqslant j \leqslant p-1\) ）。这由以下引理取 \(\mathbf{H} = \mathbf{H}_i\) , \(\mathbf{H}' = \mathbf{H}_{i+1}\) , \(\mathbf{K} = \mathbf{K}_j\) , \(\mathbf{K}' = \mathbf{K}_{j+1}\) .

引理1（Zassenhaus）。设 H 和 K 是带算子群 G 的两个稳定子群，H' 和 K' 分别是 H 和 K 的正规稳定子群；则 H'.(H ∩ K') 是 H'.(H ∩ K) 的一个正规稳定子群，K'.(K ∩ H') 是 K'.(K ∩ H) 的一个正规稳定子群，并且带算子的商群 (H'.(H ∩ K))/(H'.(H ∩ K')) 与 (K'.(K ∩ H))/(K'.(K ∩ H')) 是同构的。

由定理4， \(\mathbf{H}' \cap \mathbf{K} = \mathbf{H}' \cap (\mathbf{H} \cap \mathbf{K})\) 是 \(\mathbf{H} \cap \mathbf{K}\) ；类似地 \(\mathbf{K}' \cap \mathbf{H}\) 是 \(\mathbf{K} \cap \mathbf{H}\) 的一个正规稳定子群；因此（第6节，推论2） \((\mathbf{H}' \cap \mathbf{K})(\mathbf{K}' \cap \mathbf{H})\) 是 \(\mathbf{H} \cap \mathbf{K}\) 。将定理4应用于群 \(\mathbf{H}\) ,

\[
\mathbf {H} ^ {\prime}. (\mathbf {H} ^ {\prime} \cap \mathbf {K}). (\mathbf {K} ^ {\prime} \cap \mathbf {H}) = \mathbf {H} ^ {\prime}. (\mathbf {H} \cap \mathbf {K} ^ {\prime})
\]

是 \(\mathbf{H}'\) . (H ∩ K) 和商群

\[
(\mathbf {H} ^ {\prime}. (\mathbf {H} \cap \mathbf {K})) / (\mathbf {H} ^ {\prime}. (\mathbf {H} \cap \mathbf {K} ^ {\prime}))
\]

（使用复合

\[
(\mathbf {H} \cap \mathbf {K}) / ((\mathbf {H} ^ {\prime} \cap \mathbf {K}). (\mathbf {K} ^ {\prime} \cap \mathbf {H})).
\]

在最后一个商中，H 与 H' 一方面，K 与 K' 另一方面是对称出现的；将它们互换即可得到所述结果。

定义10. 带算子群 G 的一个 Jordan-Hölder 列是一个严格递减的分解列， \(\Sigma\) 使得不存在不同于 \(\Sigma\) 且比其更细的严格递减分解列， \(\Sigma\) .

命题9. 对于 G 的一个分解列，它是 Jordan-Hölder 列的充分必要条件是：该列的所有商都是单的。

一个分解列是严格递减的，当且仅当它的相继商中没有一个是单位元。若一个严格递减的合成列 \(\Sigma\) 不是 Jordan-Hölder 列，则存在一个严格递减的合成列， \(\Sigma'\) 它比 \(\Sigma\) 更细且不同于 \(\Sigma\) 。因此存在两个相继项 \(G_i\) , \(G_{i+1}\) 于 \(\Sigma\) 它们在 \(\Sigma'\) 中不是相继的；设 H 是紧随 \(G_i\) 在 \(\Sigma'\) 之后的第一个项；H 是 \(G_i\) 的一个正规稳定子群，且包含 \(G_{i+1}\) 且与后者不同；因此 H/ \(G_{i+1}\) 是 \(G_i/G_{i+1}\) ，与后者不同且不同于单位元；因此 \(G_i/G_{i+1}\) 不是单群。反之，如果 \(\Sigma\) 是一个严格递减的合成列，其一个商 \(G_i/G_{i+1}\) 不是单的，这个商群包含一个不同于自身的正规稳定子群。 \(\{e\}\) ，其在 \(G_i\) 是 G 的一个正规稳定子群 H， \(G_i\) ，不同于 \(G_i\) 和 \(G_{i+1}\) (定理4)；只需在H之间插入 \(G_i\) 和 \(G_{i+1}\) 为了得到一个严格递减的合成列，且不同于 \(\Sigma\) 且比其更细的严格递减分解列， \(\Sigma\) .

定理6（Jordan-Hölder）：一个带算子群的任意两条Jordan-Hölder序列是等价的。

设 \(\Sigma_{1}, \Sigma_{2}\) 设G是一个带算子群的两条Jordan-Hölder序列；通过应用定理5，两条等价的合成序列 \(\Sigma_{1}^{\prime}, \Sigma_{2}^{\prime}\) 分别比 \(\Sigma_{1}\) 和 \(\Sigma_{2}\) 更细；由于后者是Jordan-Hölder序列， \(\Sigma_{1}^{\prime}\) 与……相同 \(\Sigma_{1}\) 或者是由它通过重复某些项而得到的；其商的序列 \(\Sigma_{1}^{\prime}\) 是由……的推导而来 \(\Sigma_{1}\) 通过插入若干与群同构的项 \(\{e\}\) ；因为 \(\Sigma_{1}\) 严格递减，因此……的商序列 \(\Sigma_{1}\) 是由 \(\Sigma_{1}^{\prime}\) 的推导而来，通过抑制后者中所有同构于 \(\{e\}\) 类似地，对于 \(\Sigma_{2}\) 和 \(\Sigma_{2}^{\prime}\) 。由于商序列 \(\Sigma_{1}^{\prime}\) 和 \(\Sigma_{2}^{\prime}\) （在同构意义下）仅在各项的顺序上有所不同，这一点同样适用于 \(\Sigma_{1}\) 和 \(\Sigma_{2}\) 定理得证。

推论。设G是一个具有算子且存在若尔当-赫尔德列的群。若 \(\Sigma\) 对于G的任何严格递减合成列，存在一个比其更细的Jordan-Hölder列。 \(\Sigma\) .

设 \(\Sigma_0\) 是 \(G\) ；根据定理5，存在两个等价的合成序列， \(\Sigma'\) 和 \(\Sigma_0'\) 分别比……更细 \(\Sigma\) 和 \(\Sigma_0\) 定理6的论证表明，通过从 \(\Sigma'\) 重复，一个序列 \(\Sigma''\) 得到，这与……等价。 \(\Sigma_0\) 因此它是一个Jordan-Hölder序列，因为其所有商都是单的（命题9）。 \(\Sigma\) 是严格递减的， \(\Sigma''\) 比……更细 \(\Sigma\) ，由此得出推论。

注：带算子的群并不总是具有Jordan-Hölder列；一个例子由加法群给出。 \(\mathbf{Z}\) 有理整数的序列：该序列 \((2^{n}.\mathbf{Z})_{n\geqslant 0}\) 是 \(\mathbf{Z}\) ；对于所有 \(p\) 的（正规）子群的严格递减无穷序列，其前 \(p\) 项与该群一起构成 \(\{0\}\) 一个严格递减的合成列；如果Z存在一个Jordan-Hölder列，那么它至少会有 \(p + 1\) 由定理6的推论，各项；荒谬，因为p是任意的。

另一方面，在每一个带算子的有限群G中都存在一个Jordan-Hölder合成列：若 \(G \neq \{e\}\) 在G的异于G的正规稳定子群中，设 \(H_{1}\) 是极大子群；类似地定义 \(H_{n+1}\) 通过归纳法，作为 \(H_{n}\) 不同于 \(H_{n}\) 的正规子群集合中的极大元素，当 \(H_{n} \neq \{e\}\) ; 这些阶的序列 \(H_{n}\) 严格递减，因此存在一个n使得 \(H_{n} = \{e\}\) 由G以及 \(H_{i} (1 \leqslant i \leqslant n)\) 按其构造，这是一个Jordan-Hölder序列。

定义11. 设G是一个带算子的群；G的长度是使得存在G的严格递减合成列（G \(_{i}\) ) \(_{0 \leq i \leq n}\) .

若G具有Jordan-Hölder列，则G的长度为该列的相继商的数量，这由定理6的推论得出。若G不具有Jordan-Hölder列，则其长度为无穷；根据命题9，对于G的每个严格递减列，都存在一个严格更细的严格递减列。由单位元构成的群是唯一长度为0的带算子群。一个带算子群是单群当且仅当其长度为1。

设G是一个带算子的群，H是G的正规稳定子群，K是商群G/H，且 \(\pi: G \rightarrow K\) 典范同态。设

\[
\Sigma^ {\prime} = \left(\mathrm{H} _ {i}\right) _ {0 \leqslant i \leqslant n}
\]

是 H 的一个分解列，且 \(\Sigma'' = (\mathbf{K}_j)_{0 \leqslant j \leqslant p}\) 设K的一个合成列，记作 \(G_i = \pi^{-1}(\mathbf{K}_i)\) 对于 \(0 \leqslant i \leqslant p\) 和 \(G_i = H_{i-p}\) 对于 \(p \leqslant i \leqslant n + p\) ，一个合成列 \(\Sigma = (G_i)_{0 \leqslant i \leqslant n+p}\) 就得到了。其商序列 \(\Sigma\) 是通过将 \(\Sigma''\) 的商序列与 \(\Sigma'\) 的一个正规子群。若 \(\Sigma'\) 和 \(\Sigma''\) 的商序列并列而得到的， \(\Sigma\) 是G的Jordan-Hölder序列，由命题9可知。若H或K具有任意长度的合成序列，则G亦然。我们已证明：

命题10. 设G是一个带算子的群，H是G的一个正规稳定子群。G的长度等于H与G/H的长度之和。

推论. 设G是一个带算子的群，且 \((\mathrm{G}_{i})_{0\leqslant i\leqslant n}\) G 的一个合成列。G 的长度是 \(G_{i}/G_{i+1}, 0 \leqslant i \leqslant n - 1\) .

若G与G'是同构的带算子群，且G具有Jordan-Hölder列，则G'亦然，且G与G'的Jordan-Hölder列等价。然而，不同构的群也可能具有等价的Jordan-Hölder列；例如Z/4Z与(Z/2Z) × (Z/2Z)即如此，参见第10号。

